\addcontentsline{toc}{subsubsection}{Solution 44. The Hidden Golden Exponent}

\begin{dlsummary}[title={\bfseries Solution 44. The Hidden Golden Exponent},breakable]

\begin{proof}
The unique solution is
\[
\boxed{
f(x)=x^\varphi,\qquad
\varphi=\frac{1+\sqrt5}{2}.
}
\]

\paragraph{1. Extract superadditivity.}

Multiplicativity and positivity give \(f(1)=1\). The inequality with \(y=1\) becomes
\[
f(x+1)\ge f(x)+1.
\]
For \(u,v>0\), substitute \(x=u/v\), then multiply by \(f(v)\). Multiplicativity yields
\[
\boxed{f(u+v)\ge f(u)+f(v).} \tag{1}
\]
In particular, \(f\) is strictly increasing: if \(b>a>0\), then
\[
f(b)\ge f(a)+f(b-a)>f(a).
\]

\paragraph{2. Show that \(f\) is a power function.}

Let
\[
p=\frac{\log f(2)}{\log2}.
\]
Since \(f(2)\ge2\), we have \(p\ge1\).

Multiplicativity gives, for every rational \(r\),
\[
f(2^r)=f(2)^r=2^{pr}.
\]
For any \(t>0\), choose rational numbers \(r_n,s_n\) such that
\[
r_n<\log_2t<s_n,
\qquad r_n,s_n\longrightarrow\log_2t.
\]
Monotonicity implies
\[
2^{pr_n}\le f(t)\le2^{ps_n}.
\]
Thus
\[
\boxed{f(t)=t^p\qquad(t>0).} \tag{2}
\]
This establishes continuity as a consequence, rather than assuming it.

\paragraph{3. Pin down the exponent.}

Substituting (2) into the original inequality gives
\[
(x+y^p)^p\ge x^p+y^{p+1}.
\]
For fixed \(y>0\), let \(x\to0^+\). We obtain
\[
y^{p^2}\ge y^{p+1}\qquad(y>0).
\]
Choosing \(y>1\) gives \(p^2\ge p+1\), while choosing \(0<y<1\) gives the reverse inequality. Hence
\[
p^2=p+1.
\]
Since \(p\ge1\),
\[
p=\varphi=\frac{1+\sqrt5}{2}.
\]

\paragraph{4. Verify the solution.}

Multiplicativity is immediate. For any \(a,b>0\) and \(p\ge1\),
\[
(a+b)^p
=a(a+b)^{p-1}+b(a+b)^{p-1}
\ge a^p+b^p.
\]
Taking \(a=x\), \(b=y^\varphi\), and using \(\varphi^2=\varphi+1\), we get
\[
(x+y^\varphi)^\varphi
\ge x^\varphi+y^{\varphi^2}
=x^\varphi+y^{\varphi+1},
\]
exactly the required inequality.
\end{proof}
\end{dlsummary}

\begin{takeaways}
\begin{itemize}
\item \textbf{A unit shift can control arbitrary sums:} Multiplicativity converts \(f(x+1)\ge f(x)+1\) into superadditivity.
\item \textbf{Positivity supplies regularity:} Superadditivity forces monotonicity, eliminating nonstandard multiplicative functions.
\item \textbf{Use both sides of \(1\):} A power comparison valid for every positive input often forces equality of exponents.
\item \textbf{Composition multiplies exponents:} The competition between \(f(f(y))\) and \(y f(y)\) produces \(p^2=p+1\).
\end{itemize}
\end{takeaways}
